\documentclass[a4paper,12pt]{article}

\usepackage[utf8]{inputenc} % encodé en utf-8
\usepackage[T1]{fontenc} % compatible avec les accents

\usepackage[french]{babel} % rédigé en français
\usepackage[hyphens]{url} % formatte les liens en autorisant la césure au niveau des traits d'union
\usepackage[urlcolor=black,colorlinks=true,linkcolor=black,citecolor=black]{hyperref} % liens cliquables mais non colorés

\usepackage[top=25mm, bottom=25mm, left=35mm, right=25mm]{geometry} % définit les marges
\usepackage{array} % permet l'alignement vertical dans les tableaux
\usepackage{setspace} % permet de régler l'interligne
\onehalfspacing % définit l'interligne 1,5 par défaut

\usepackage{graphicx} % gestion des images
\graphicspath{ {../images/}}
\usepackage{array} % gestion des tableaux
\usepackage{csquotes} % gestion des guillemets

\usepackage{fourier} % utilise une autre police que celle par défaut (Computer Modern) fourier - venturis - libertine - tgtermes - kpfonts - lmodern - txfonts - newtxtext,newtxmath - mathptmx - stix 

\usepackage[backend = biber, style = authoryear, useprefix = true, maxbibnames = 99]{biblatex}


%\addbibresource{}

\DefineBibliographyExtras{french}{%
\protected\def\mkbibnamefamily#1{%
\textsc{\textnohyphenation{#1}}}
}%

\DefineBibliographyStrings{french}{%
    page  = {\lowercase{p.}}, % for single page number
    pages = {\lowercase{pp.}}, % for multiple page numbers
}

\DefineBibliographyExtras{french}{\renewcommand*{\bibrangedash}{\textendash}} % Pour rendre les "--" correctement

% \renewcommand*{\mkibid}{\emph} %ibid en italique  if authoryear-ibid

\usepackage{float}

\setcounter{secnumdepth}{3} %Ne compte pas plus loin que 1.1.1
\setcounter{tocdepth}{3}

\usepackage{enumitem} %% utilisé pour la présentation de fausse table des matières dans la table des matières commentée

\newlist{legal}{enumerate}{10}
\setlist[legal]{label*=\arabic*.}

\usepackage{ragged2e} % Permet de placer plus facilement les tableaux

\newcolumntype{R}[1]{>{\RaggedRight}p{#1}}

\title{Projet STIC-B505 : Budgetsquirrel, une application de gestion de budget} 
\author{Milena Albu \and Marie Huysman} % vos prénom et nom
\date{} % pas de date

%\usepackage{fancyhdr}
%\pagestyle{fancy} % Turn on the style
%\fancyhf{}
%\fancyhead[L]{\footnotesize \nouppercase \leftmark}
%\fancyhead[R]{\footnotesize \nouppercase \rightmark}
%\fancyfoot[C]{\thepage}
%
%\renewcommand{\chaptermark}[1]{%
%\markboth{#1}{}}
%
%    \renewcommand{\sectionmark}[1]{\markright{\thesection \ #1}{}}

\usepackage{longtable}


\begin{document} % début du corps du texte
\sloppy

\begin{titlepage} % style spécial pour la page de garde
\singlespacing % remet l'interligne normal pour cette page

%%alternative à la bannière image
%\begin{center} 
%\huge{\textsc{Université Libre de Bruxelles}}\\
%\vspace{.5cm}
%\LARGE{Faculté de Lettres, Traduction et Communication}
%\end{center}
%% Bannière image
\begin{figure}[h]
\noindent
\makebox[\textwidth]{\includegraphics[scale=0.6]{banner.png}}
\end{figure}

\vfill % laisse LaTeX gérer l'espacement vertical

%VOIR SI LE TITRE CONVIENT

\begin{center}
\LARGE{Conception et gestion de banques de données : Rapport de projet}\\ % ajoutez votre titre ici
\Large{ Budget Squirrel : Application de gestion de budget}
\vspace{.5cm}
\end{center}

\vfill

\begin{tabular}{b{5.5cm}b{7.5cm}} % taille et alignement des cellules du tableau
 \\ Milena \textsc{Albu} \\ Marie \textsc{Huysman} \\ & Rapport de projet écrit dans le cadre du cours STIC-B505 : Conception et gestion de banques de données \\
\end{tabular}

\vfill

\begin{center}
Année académique 2020--2021 % vérifiez que l'année est correcte !
\end{center}

\end{titlepage}

\tableofcontents

\thispagestyle{plain}

\newpage

\section{Introduction}

Si besoin de citer : \parencite[2--4]{SOURCE}
Notes de bas de page \footnote{informations suppl.}

Dans le cadre de ce projet, nous avons choisi de développer une application permettant la gestion de budget à une échelle personnelle.
Après une brève présentation du domaine d'application, nous présenterons le schéma conceptuel, ainsi que sa traduction en schéma relationnel.
Ensuite, nous décrirons les différents aspects des scripts SQL que nous avons mis en place [...] .
Enfin, nous exposerons brièvement l'application Web que nous avons développée. 
Vous trouverez, en annexe à ce rapport, deux scripts sql : \verb\initdb.sql\, qui contient tous les scripts permettant la création de  la base de donnée en elle-même, et \verb\data.sql\, qui peuple la base de données de quelques informations, permettant ainsi de parcourir les différentes pages de l'application.

\subsection{Domaine d'application}

Utilisateurs cibles, cas d'utilisation du système

\newpage

\section{Schéma conceptuel}

Schéma conceptuel documenté, accompagné d'une description et justification des choix

\newpage 

\section{Schéma logique relationnel}

Schéma logique relationnel traduisant le schéma conceptuel, avec une explication du choix des clés de traductions (héritage, associations, contraintes...

\newpage

\section{Code SQL}

\subsection{Description du code de création de la base de données}

Description du code sql de création de la db 
!! La création de la db doit : créer des tables , avoir des contraintes de clé primaire, unique et externe, avoir des contraintes sous forme de prédicat de table et de colonne, avoir des procédures stockées et des déclencheurs

\subsection{Éléments de SQL avancé}

\subsubsection{Contraintes garanties par des check}

\subsubsection{Requêtes de consultation}

\subsubsection{Requêtes de mise à jour}


\subsubsection{Vues utilisées}

Ce qu'on utilise pour l'historique et les différentes statistiques

\subsubsection{Déclencheurs}

Triggers à utiliser, sans doute, pour le calcul des bilans, et autres 

\subsubsection{Droits d'accès}

Voir si les droits d'accès peuvent être utilisés comme dans le rapport publié par Servais, pour l'instant ça reste un accès root défini par une session > problématique?



\newpage
\section{Description de l'application Web}

L'application web que nous avons développé a pour objectif de permettre aux utilisateurs de s'inscrire, d'enregistrer leurs transactions, de consulter leur historique par mois, ainsi que des statistiques sur la répartition de leurs dépenses, de leur revenu, et un bilan total de leur budget.

À l'ouverture de l'application (\verb\localhost/budgetsquirrel/index.html\ ou simplement \verb\localhost/budgetsquirrel/\), l'utilisateur a deux possibilités : se connecter, ou s'inscrire.

S'il choisit de s'inscrire, il est redirigé vers un formulaire d'inscription où il doit rentrer son nom, son prénom, son NISS, et choisir une photo de profil. Si un utilisateur est déjà inscrit sous le NISS entré, il en est notifié (et a la possibilité de se rediriger vers la page de connexion). Si le NISS n'est pas encore utilisé par un utilisateur, l'enregistrement se fait, l'utilisateur est notifié du succès de ce dernier, et peut se rediriger vers la page de connexion.

La page de connexion, par souci de simplification, est une simple sélection de profil : l'utilisateur a accès à la liste des utilisateurs, sélectionne son profil (nom et prénom) dans la liste, et se connecte (à ce niveau, le niss de l'utilisateur est récupéré par une méthode POST, enregistré par une méthode SESSION et appelé grâce à \verb\session_start();\ aux différentes pages de l'application).

\textbf{NOTE : A AJOUTER DANS LA DB : UN ACCES UTILISATEUR A UTILISER PARTOUT DANS L'APP vs un autre accès admin - DB à utiliser slmt au niveau de la db (qui peut par ex. supprimer des utilisateurs et ajouter de nouvelles catégories) comme ça , on exploite aussi les droits d'accès qui sont demandés par le projet et on évite l'accès root, en délimitant ce que l'utilisateur peut consulter et modifier}

Une fois qu'il a sélectionné son profil et choisi de se connecter, l'utilisateur est redirigé vers une page d'accueil. La page d'accueil présente brièvement les différentes pages de l'application. Depuis la barre de navigation, présente sur toutes les pages (sauf celles ou l'utilisateur n'est pas connecté), l'utilisateur a accès : à la page d'accueil, à l'historique, à la page d'enregistrement, à la page de statistiques, et à sa page personnelle de profil.

La page personnelle de profil permet à l'utilisateur de visualiser ses informations : nom, prénom, NISS, date de naissance. Il peut également y ajouter et y supprimer ses cartes\footnote{Cette suppression est une suppression logique du côté de la base de données : les cartes possèdent une colonne is\_deleted , par défaut la valeur de la colonne est à 0, et si l'utilisateur \og supprime\fg sa carte, la valeur passe à 1, et la carte est considérée comme virtuellement supprimée} . Les listes de ses cartes disponibles et de ses anciennes cartes sont également visibles depuis la page de profil. Il peut également y actualiser sa photo.

La page d'enregistrement permet d'enregistrer des transactions. L'utilisateur est obligé d'entrer un montant, de sélectionner une date, une catégorie de transaction, et un type de transaction. En fonction du type de transaction effectué, il doit également ajouter des informations supplémentaires. Le payement en liquide ne demande aucune information supplémentaire, le virement bancaire demande d'introduire un destinataire ou un bénéficiaire, et éventuellement une communication, et un payement par carte nécessite de sélectionner la carte utilisée.

La page historique permet à l'utilisateur de visualiser les transactions effectuées : pour cela, il doit d'abord sélectionner le mois et l'année dont il veut consulter le budget. Il a alors accès à un tableau affichant le montant, la date, la catégorie, et le type de transaction effectuée, et éventuellement la carte utilisée, le destinataire/bénéficiaire, et la communication. En dessous de la table, l'utilisateur peut également voir le bilan total du mois. Enfin, il peut également choisir de supprimer une des transactions affichées.

L'écran de statistiques offre différentes informations sur les données concernant l'utilisateur. Les informations offertes sont les suivantes : un total des dépenses, un total des revenus, et un bilan, mais aussi un tableau montrant la répartition des transactions (dépenses et revenus) par mois, ainsi que le bilan par mois, pour une vue plus condensée de l'historique, un tableau montrant la répartition totale par catégorie (couplé à une description de chaque catégorie), et  enfin la répartition des transactions entre les différents types de payements (le nombre d'utilisations, le bilan des dépenses et le bilan des revenus pour chaque type de transaction).

À tout moment, l'utilisateur peut également se déconnecter, en cliquant, en haut à droite de la barre de navigation, sur "déconnexion". Les informations de sessions (donc, le NISS sous lequel l'utilisateur est connecté) sont effacées, et l'utilisateur peut choisir de retourner à l'écran d'accueil.

\newpage

\section{Développement du projet}

Quelques mots sur la logique du travail : paired programming, outils employés, défis rencontrés, ...

\newpage
\section{Conclusion}

Les points forts et faiblesses d’application, le temps nécessaire pour développer le concept avec deux personnes, les sources consultées pour la partie pratique (i.e. les slides, ou techniques vues au TP, ainsi que les conseils d’assistante, mais aussi w3schools ou bien stackoverflow).
Finalement quelques mots sur l'expérience du projet.




\newpage

%voir si bibliographie nécessaire
%\chapter*{Bibliographie} 
%\addcontentsline{toc}{chapter}{Bibliographie}
%\markboth{Bibliographie}{}
%
%\printbibliography[heading=none]
%
%\fussy

%Voir si annexes nécessaire
%\chapter*{Annexes}
%\addcontentsline{toc}{chapter}{Annexes}
%\markboth{Annexes}{}

\end{document} % fin du corps du texte